% !TEX root = ../main.tex
\chapter{Data formats and units}\label{app:data}
This appendix says how recordings and everything made from them are stored. The rule behind all of
it: bulk arrays live on disk as files, and the database holds paths and small facts
(\cref{sec:architecture-storage}). Paths are relative to the repository root. \code{DATA/},
\code{MODELS/}, \code{MATRICES/} and \code{Plots/} are not in version control; each machine
provisions its own. The sources are \code{docs/DATA_REGISTRATION.md},
\code{docs/LIBRARY_STORAGE.md}, \code{Working/database/schema.py} and \code{Working/units.py}.

\section{Recordings on disk}\label{sec:data-recordings}
A \term{recording}{recording} arrives as a raw file and is \emph{derived} into one array per
\term{channel}{channel} before the app can read it.

\begin{description}
\item[Raw files] \code{DATA/raw/}, named \code{<stem>[_fs<N>].mat} or \code{.csv}. The layouts met
  so far: a flat vector of sixteen channels end to end (\code{M2_aug_concat_fs1.mat}), a matrix of
  samples by channels, a channels-by-samples MATLAB v7.3 matrix, and one CSV column per channel.
  The \term{sampling rate}{sampling-rate} comes from the \code{_fs<N>} in the name or is supplied at
  import; a raw file whose rate cannot be known is refused.
\item[Derived channels] \code{DATA/derived/channels/<stem>/}: one file \code{CH<n>.npy} per channel,
  a one-dimensional float64 array, with \code{<n>} the zero-based channel index, plus a
  \code{manifest.json} giving \code{source_file}, \code{fs}, \code{n_channels},
  \code{n_samples_per_channel}, \code{dtype}, \code{fs_note}, \code{time_base} and \code{units}.
  Importing a raw file writes this folder (staged, then swapped in at once) and registers it.
\item[The \code{recordings} table] One row per channel, unique on \code{(source_file, channel)}:
  \code{fs}, \code{n_samples}, \code{global_offset}, \code{npy_path} (repository-relative, which is
  why the bridge runs from the repository root), \code{notes}, and, added later, registration
  provenance (\code{fs_source} read or inferred, \code{registered_at}, \code{registered_by},
  \code{warnings_json}, \code{active}), the excerpt link (\code{parent_recording_id},
  \code{parent_offset}, \code{decimation}) and the unit (\code{units}, \code{units_note}).
\end{description}

The core reads a channel by memory-mapping its \code{.npy} file and slicing the span it needs
(\code{_load_signal} in \code{Working/execution.py}); recordings are never written by a run.

\begin{plainterms}[Channel names count from one]
Files and the database count channels from zero, as programmers do: \code{CH2.npy} is stored channel
2. The app prints names counting from one, as people do, so that file is ``CH3''
(\code{channel_name} in \code{Working/discovery/channels.py}). Where it matters, the stored index is
shown beside the name.
\end{plainterms}

A \term{dataset}{dataset} is the set of channel rows that share a \code{source_file}. Its identity
(display name, species, organism id, experiment date, condition, notes) is one row of the
\code{datasets} table, keyed by source file and created the first time a field is filled. A display
name is a label only; nothing keys, joins or caches on it.

\begin{description}
\item[Excerpts] A channel that is a shorter, downsampled piece of another registered channel is an
  \emph{excerpt}, not a new recording. The registration check (\code{Working/registration/excerpt.py})
  block-averages the longer channel to the shorter one's rate and slides one along the other; a
  correlation above 0.99 that clearly beats its neighbours marks an excerpt, which is linked to its
  parent rather than registered again.
\item[Held out] \code{M4_aug_concat_fs1.mat} is the \term{held-out}{held-out} recording. Registration
  refuses it, every Explore and Analyse route refuses its rows with HTTP 423, and the executor
  refuses it unless \code{HELD_OUT_UNLOCK} is set in \code{Working/config.py}.
\end{description}

\section{The sidecar manifest}\label{sec:data-manifest}
Every registered file gets a provenance file beside it, written by \code{register}:
\code{<file>.manifest.json} next to a file, or \code{<dir>/registration.manifest.json} inside a
folder. The database row points at it (\code{registered_artifacts.manifest_path}). Its fields, from
\code{docs/DATA_REGISTRATION.md}:

\begin{description}
\item[\code{manifest_version}, \code{kind}, \code{name}, \code{path}] what the file is and where.
\item[\code{source}] the recording, its row id, channel, span, and the raw file, if any.
\item[\code{fs}, \code{fs_source}] the sampling rate and whether it was read or inferred.
\item[\code{parameters}] what the producer used (for a matrix profile: window length, config hash,
  backend).
\item[\code{producer}, \code{recipe_hash}] the recipe, script or cluster job that made the file.
\item[\code{created_at}, \code{registered_by}, \code{code_version}] when, by whom, at which commit.
\item[\code{checks_passed}, \code{warnings}, \code{sha1}, \code{row}] the checks it passed, the
  warnings raised, a fingerprint of the file, and the database row as inserted.
\end{description}

The fingerprint of a file is its SHA-1; of a folder, each file's name, size and the SHA-1 of its first
and last 4\,MB. In the bridge's sandbox mode, sidecars are written under
\code{webui/runtime/<stamp>/sidecars/} instead (\cref{sec:architecture-sandbox}).

Not to be confused with this: the per-recording \code{manifest.json} in a derived-channels folder
(above), and the run manifest a cluster job writes through \code{Pipelines/run_recipe}
(\code{Working/manifest.py}).

\section{Units}\label{sec:data-units}
\code{recordings.units} is \code{V}, \code{mV} or \code{uV}; empty means \emph{undeclared}.
\code{recordings.units_note} says where the unit came from, or why it is unknown. A unit comes from,
in order: the derived folder's manifest (\code{units}; free text is read by its first word, so
``millivolts as stored \ldots'' is \code{mV}), an override at registration, or Settings
\textrightarrow{} Datasets afterwards (\code{PUT /api/registry/recording/{id}/units}, which sets
every channel of the dataset together and writes an audit entry). An unreadable unit fails
registration's \code{units} check; an undeclared one is a warning.

\begin{center}\small
\begin{tabular}{@{}lll@{}}
\toprule
\bfseries Stored unit & \bfseries Factor to mV (\code{to_mv_factor}) & \bfseries Read by a block that prints mV\\
\midrule
\code{V} & 1000 & $\times 1000$\\
\code{mV} & 1 & $\times 1$\\
\code{uV} & 0.001 & $\times 0.001$\\
undeclared & none & read as volts, $\times 1000$, with a note saying so\\
\bottomrule
\end{tabular}
\end{center}

The core never converts stored samples. The bridge converts only what it draws
(\code{display_channel} in \code{webui/server/corpus.py}) and labels an undeclared recording ``unit
undeclared'' rather than ``mV''. A block that prints millivolts is told the unit and converts with
\code{read_as_mv}; every other block works in the stored unit (\cref{sec:architecture-units}).
\code{init_db()} fills in the unit measured on 2026-09-23 for each known file
(\code{RECORDING_UNITS_EVIDENCE} in \code{Working/units.py}) on rows with no unit and no note, so a
declared unit is never overwritten. At the time of writing, the M2 files and
\code{Mushroom_260720_0509_4hrs_CH14_fs1.mat} are recorded as volts, \code{L_LM_Jul_26_J_raw.mat}
as millivolts, and the others (including M1, M100, M101\_t, MJu26a and the synthetic
\code{Fig2A_dt0p1.csv}) as undeclared, each with its reason. Recordings registered since then by
the derivation scripts declare their own unit in their manifest.

\begin{watchout}[A threshold is in the stored unit]
A block that is not unit-aware compares its threshold with the samples as stored. The same number is
a thousand times stricter on a millivolts recording than on a volts one. Each block's entry in
\cref{app:blocks} says whether it is unit-aware.
\end{watchout}

\section{The database}\label{sec:data-database}
One SQLite file, \code{DATA/db/annotations.sqlite}, plain SQL, its whole schema in
\code{Working/database/schema.py}, created and migrated by \code{init_db()}
(\cref{sec:architecture-storage}). The tables a researcher meets most:

\begin{description}
\item[\code{recordings}, \code{datasets}] the channels and their datasets (above).
\item[\code{configs}] one row per distinct recipe: \code{config_hash} (the
  \term{recipe}{recipe}'s eight-character short hash, unique) and \code{config_json} (the recipe as
  canonical JSON).
\item[\code{runs}] one row per execution of a recipe over a recording and span: \code{config_id},
  \code{recording_id}, \code{span_start}, \code{span_end}, \code{status}, \code{started_at},
  \code{finished_at}, \code{duration_s}, \code{step_timings_json}, \code{error_text} (a full
  traceback on failure), \code{current_step}, a fan-out group, a surrogate pairing, an optional name,
  and \code{superseded_at} for a discarded run.
\item[\code{detections}] machine output: \code{run_id}, \code{start_idx}, \code{end_idx} (sample
  indices in the whole channel), \code{score}, \code{meta_json}. Written only by runs.
\item[\code{artifacts}] files a run produced: \code{run_id}, \code{kind} (plot, encoding, model,
  csv, other), \code{path} (repository-relative).
\item[\code{annotations}] human judgements of a span: \code{recording_id}, \code{start_idx},
  \code{end_idx}, \code{verdict} (one of the five, enforced by a \code{CHECK}), tag, note, the scale
  viewed, source, time; deleting sets \code{deleted_at} rather than removing the row.
\item[\code{adjudications}] a human verdict on one machine detection (\code{detection_id},
  unique), kept apart from both \code{detections} and \code{annotations}
  (\cref{sec:core-ideas-detections}).
\item[\code{tag_vocabulary}] the controlled tags, linked to annotations, adjudications and motifs.
\item[\code{motif_entry}, \code{motif_member}, \code{motif_edge}, \code{motif_member_revision},
  \code{groupings}, \code{grouping_assignments}, \code{hand_edits}, \code{sequences},
  \code{sequence_members}] the Library: shapes, their occurrences, relationships, revisions and
  saved groupings. An entry's identity is a hash of its z-normalised, resampled waveform
  (\code{docs/LIBRARY_STORAGE.md} \S2).
\item[\code{window_sets} and its member tables] saved training-window sets; their bulk bounds are on
  disk.
\item[\code{review_queues}, \code{window_verdicts}, \code{review_audit}, \code{blind_test_labels}]
  Review's queues and the verdicts given in them.
\item[\code{discovery_sessions}, \code{discovery_runs}, \code{run_groups}] Discovery's scopes and the
  runs fanned out from one recipe.
\item[\code{templates}] saved chains with recording and span removed.
\item[\code{jobs}] the bridge's long work: kind, status, times, progress, error, linked run.
\item[\code{step_artifacts}] the step cache's index (below).
\item[\code{registered_artifacts}, \code{encodings}] registered files other than recordings: path,
  manifest path, recording, channel, span, settings, producer, SHA-1, checks, warnings,
  \code{active}.
\item[\code{settings}] project settings: page, key, JSON value, time, actor (\cref{ch:settings}).
\item[\code{audit_log}] append-only record of changes: time, kind, what, where, route, actor, detail.
\end{description}

\begin{underthehood}[Which code writes which table]
Machine tables and human tables are written through two separate doors, \code{write_machine} and
\code{write_human} in \code{webui/server/writes.py}, each refusing the other's tables. The executor
writes runs, detections and artifacts through \code{Working/database/runs.py}. Library writes go
through \code{webui/server/library.py}; \code{docs/LIBRARY_STORAGE.md} \S3.5 lists the door for each
Library table.
\end{underthehood}

\section{Derived files and caches}\label{sec:data-derived}
Everything below is made by the app from the recordings, and can be made again. Default locations,
in \code{--project} mode:

\begin{center}\small
\begin{tabular}{@{}>{\raggedright\arraybackslash}p{0.42\linewidth}>{\raggedright\arraybackslash}p{0.52\linewidth}@{}}
\toprule
\bfseries Where & \bfseries What\\
\midrule
\code{DATA/derived/step_cache/<prefix-hash>/<step>/} & cached typed output of a slow step, indexed by \code{step_artifacts} (\code{STEP_CACHE_ROOT})\\
\code{Results/Detection/matrix_profile/} & matrix profiles, \code{mp_v2_<stem>_CH<n>_WIN<len>min[_span<a>-<b>].npz}\\
\code{Results/Preprocessing/window_matrix/} & window matrices, \code{wm_v1_<stem>_CH<n>_WIN<len>min_STEP<pct>pct.npz}; older ones in \code{MATRICES/}\\
\code{DATA/derived/encodings/} & cached encodings, \code{enc_<type>_<stem>_CH<n>_<hash8>.npz}\\
\code{DATA/derived/window_sets/<name>/} & window sets: \code{windows.npz} and \code{manifest.json}\\
\code{DATA/derived/models/}, \code{MODELS/} & trained models (\code{.joblib}, \code{.pth})\\
\code{DATA/derived/groupings/}, \code{shape_trees/}, \code{trace_shapes/} & clustering labels, shape trees, trace-shape vectors\\
\code{DATA/derived/drop_motifs/<run>/} & drop-detector event stores: \code{events.csv}, \code{snippets.npz}, \code{manifest.json}\\
\code{DATA/db/*_calibration.json} & per-machine cost calibrations (matrix profile, window matrix, other blocks)\\
\code{DATA/db/backups/} & database backups; the bridge writes one on every \code{--project} start and keeps ten\\
\code{HPC/} & cluster jobs the site wrote, and results copied back (\cref{ch:jobs})\\
\code{webui/runtime/<stamp>/} & a sandbox: its database copy, caches, results, sidecars, server log\\
\bottomrule
\end{tabular}
\end{center}

The step cache has no eviction: an entry stays until it is cleared (Settings \textrightarrow{}
Analysis defaults shows its size and can clear it). A missing or unreadable entry is recomputed and
overwritten. In sandbox mode the step cache, the results folders of the blocks that write files, the
classifier's model folder and all three cost-calibration files (matrix profile, block cost, window
matrix) are redirected into the runtime folder, and the bridge refuses to start if any of them
points outside it (\cref{sec:architecture-sandbox}).

\begin{watchout}[Do not move a file the database points to]
A database row holds a path. Moving or renaming a derived file, or the folder of a recording, leaves
the row pointing at nothing. Register the file again from its new place, and unregister the old row,
rather than moving it.
\end{watchout}
